\documentclass[11pt]{article}  

%%%%%%%%%%%%%% PREÁMBULO %%%%%%%%%%%%
\title{Propuesta de Proyecto}
\usepackage[spanish]{babel}
\usepackage[utf8]{inputenc}
\usepackage[T1]{fontenc}     % fuentes con soporte de acentos mejorado
\usepackage{listingsutf8}
\usepackage{textcomp}      % para algunos símbolos de texto
\usepackage{upquote}    
\usepackage{amsmath}
\usepackage{amssymb} 
\usepackage{graphicx}
\usepackage{color}
\usepackage{enumitem}
\usepackage{subfigure} 
\usepackage{float} 
\usepackage{capt-of} 
\usepackage{sidecap} 
\sidecaptionvpos{figure}{c} 
\usepackage{caption} 
\usepackage{commath}
\usepackage{graphicx, amsmath, amsthm, latexsym, amssymb, amsfonts, epsfig, float, enumerate, color, listings, graphicx, fancyhdr}
\usepackage{cancel} 
\usepackage[table,xcdraw]{xcolor}
\usepackage{anysize} 
\usepackage{listings}
\usepackage{xcolor}  % Para colores
\usepackage{booktabs}
\usepackage{pgfgantt}
\usepackage{pdflscape}

\marginsize{2cm}{2cm}{2cm}{2cm} % Izquierda, derecha, arriba, abajo
\usepackage{appendix}
\usepackage[style=ieee, backend=biber]{biblatex}
\renewcommand{\appendixname}{Apéndices}
\renewcommand{\appendixtocname}{Apéndices}
\renewcommand{\appendixpagename}{Apéndices} 
\usepackage[colorlinks=true,plainpages=true,citecolor=blue,linkcolor=black]{hyperref}
\usepackage{fancyhdr} 
\usepackage{url}
\usepackage{longtable}

\pagestyle{fancy}
\fancyhf{}
\fancyhead[L]{\footnotesize Ingeniería de Software para Sistemas Inteligentes} 
\fancyfoot[C]{\thepage} 
\renewcommand{\footrulewidth}{0.4pt}
\usepackage[utf8]{inputenc}

\definecolor{dkgreen}{rgb}{0,0.6,0} 
\definecolor{gray}{rgb}{0.5,0.5,0.5} 
\definecolor{verde}{rgb}{0.18,0.549,0.2}
\definecolor{rojo}{rgb}{1, 0.302, 0.431}
\definecolor{black}{rgb}{0, 0, 0}
\definecolor{blue}{rgb}{0, 0.4648, 0.71}
\definecolor{charcoal}{RGB}{54,69,79}

\definecolor{codegreen}{rgb}{0,0.6,0}          % Verde para comentarios
\definecolor{codegray}{rgb}{0.5,0.5,0.5}       % Gris para números de línea
\definecolor{codepurple}{rgb}{0.58,0,0.82}     % Púrpura para strings
\definecolor{backcolour}{rgb}{0.98,0.98,0.99}  % Color de fondo
\definecolor{keywordcolor}{rgb}{0,0.59,0.61}   % Color para palabras clave (class, def, etc.)
\definecolor{titlecolor}{rgb}{0.53,0,0}        % Color para títulos de clases y funciones
\definecolor{stringcolor}{rgb}{0,0.36,0.37}    % Color para strings
\definecolor{commentcolor}{rgb}{0.58,0.65,0.65} % Color para comentarios

% Definir el estilo para el código
\lstdefinestyle{mystyle}{
    backgroundcolor=\color{backcolour},   
    commentstyle=\color{commentcolor},
    keywordstyle=\color{keywordcolor},
    numberstyle=\tiny\color{codegray},
    stringstyle=\color{stringcolor},
    basicstyle=\ttfamily\footnotesize,
    breakatwhitespace=false,         
    breaklines=true,                 
    captionpos=b,                    
    keepspaces=true,                 
    numbers=left,                    
    numbersep=5pt,                  
    showspaces=false,                
    showstringspaces=false,
    showtabs=false,                  
    tabsize=2,
    emph={class,def,__init__,return,import}, % Palabras clave adicionales
    emphstyle=\color{titlecolor},            % Estilo para palabras clave adicionales
    morekeywords={class,def,__init__,return,import}, % Añadir palabras clave
     % ahora añadimos los símbolos faltantes
}

\lstset{style=mystyle}
\renewcommand{\arraystretch}{1.3} % Aumenta el espacio entre filas en tablas

% --- Archivo de bibliografía -----------------
\addbibresource{repbib.bib}


\title{Requerimientos Funcionales y No Funcionales}

%%%%%%%%%%%%%%%%%%%% TERMINA PREÁMBULO %%%%%%%%%%%%

\begin{document}

%%%%%%%%%%%%%%%%%%%% PORTADA %%%%%%%%%%%%%%%%%%%%%%%

\begin{titlepage}
  \thispagestyle{empty}
  \pagecolor{white}
  
  % Encabezado con logos
  \begin{center}
    \vspace{1cm}
    
    % Logos en la parte superior
   
      \centering
          \vspace{2cm}
      \includegraphics[height=2.35cm]{img/logos.png}
    
    \vspace{3cm}
    
    % Tipo de trabajo
    {\normalsize\color{charcoal}\textsc{Especificación de requisitos}}\\[0.5cm]
    
    % Título principal
    {\huge\bfseries\color{black}Documento de especificación de requisitos de software}\\[0.5cm]

    
    \vspace{1.5cm}
    
    % Información de autores
    {\small\color{charcoal}\textsc{Realizada por}}\\[0.3cm]
    
    {\normalsize\color{black}Robles Guzmán Naomi Isabel}\\[0.2cm]
    {\normalsize\color{black}Ugalde Téllez Aarón}\\[0.5cm]

    
    \vspace{0.8cm}
    
    % Información académica
    {\small\color{charcoal}\textsc{Para la materia de}}\\[0.2cm]
    {\normalsize\color{blue}Ingeniería de Software para Sistemas Inteligentes}\\[0.5cm]
    
    {\small\color{charcoal}\textsc{Impartida por}}\\[0.2cm]
    {\normalsize\color{black}Chadwick Carreto Arellano}\\[0.5cm]
    
    {\small \color{charcoal}\textsc{Grupo}}\\[0.3cm]
    {\normalsize\color{black}6BM1}\\[0.8cm]
    
    \vfill
    
    % Línea decorativa y fecha
    {\color{blue}\rule{0.4\textwidth}{1pt}}\\[0.3cm]
    {\large\color{charcoal}15 de septiembre de 2025}
    
    \vspace{1cm}
    
  \end{center}
\end{titlepage}                                                     


%%%%%%%%%%%%%%%%% CUERPO %%%%%%%%%%%%%%%%%%%%%%%%

\section{Introducción}
Él presente \textit{software} plantea una aplicación web responsiva de gestión de tareas académicas y profesionales que integra
 un motor de recomendaciones educativas basado en análisis de palabras clave. Ofreciendo una conexión directa hentre recursos educativos y 
 gestión de actividades, alineando planificación y ejecución en un mismo entorno.\\\\
 Los requisitos de \textit{software} establecen las especificaciones detalladas que definen que debe hacer un sistema, cómo debe 
 comportarse y bajo qué condiciones debe operar \autocite{Requirements2024}. Este documento  tiene como propósito establecer de
 manera precisa y completa todas las funcionalidades, características y restricciones que deberá cumplir el sistema de gestión de
 tareas académicas y profesionales, garantizando así el desarrollo de una solución que satisfaga plenamente las necesidades identificadas. 

\section{Requerimientos Funcionales}
\begin{table}[H]
    \centering
    \begin{tabular}{|p{2cm}|p{13cm}|}
    \hline
    \multicolumn{2}{|c|}{\textbf{1. Cuenta e Inicio de Sesión}} \\
    \hline
    \textbf{Clave} & \textbf{Descripción } \\
    \hline
    RF1.1 & El sistema permite al usuario crear una cuenta ingresando correo electrónico, contraseña,
     país de origen y último grado de estudios. \\
     \hline
    RF1.2 & El sistema permite al usuario iniciar sesión con su correo electrónico y contraseña registradas. \\
    \hline
    RF1.3 & El sistema permite al usuario recuperar su contraseña mediante un correo electrónico de recuperación.\\
    \hline
    RF1.4 & El sistema permite al usuario ver y actualizar su perfil, incluyendo nombre, foto de perfil, país de origen y último grado de estudios. \\
    \hline
    RF1.5 & El sistema permite al usuario terminar su sesión.\\
    \hline
    \multicolumn{2}{|c|}{\textbf{2. Registro de Tareas}} \\
    \hline
    RF2.1 & Dado un usuario autenticado, al dar cloc en "Nueva tarea", el sistema abre un modal donde el usuario puede ingresar el título, descripción
     y fecha de vencimiento de la tarea. \\
    \hline
    RF2.2 & Tras ingresar los datos de la tarea, al dar clic en "Guardar", el sistema almacena la tarea en la base de datos con la etiqueta de "Todo".\\
    \hline
    RF2.3 & Al registrar una tarea además de especificar la fecha límite de la tarea, el usuario puede especificar el plazo en el que desea realizarla, es decir,
     los días en el que desea trabajar en la tarea. \\
    \hline
    RF2.4 & Dado un usuario con tareas registradas en el mes, cuando abre la vista mensual, entonces cada día muestra una etiqueta/contador con las tareas 
    y al hacer clic se ve el detalle y acceso para editar. \\
    \hline
    RF2.5 & El sistema permite al usuario modificar tareas. \\
    \hline
    RF2.6 & El sistema permite al usuario eliminar tareas. \\
    \hline
    \end{tabular}
    \label{tab:requisitos-funcionales-1}
\end{table}

%%%%
\begin{table}[H]
    \centering
    \begin{tabular}{|p{2cm}|p{13cm}|}
    \hline
    \multicolumn{2}{|c|}{\textbf{2. Registro de Tarea}} \\
    \hline
    \textbf{Clave} & \textbf{Descripción } \\
    \hline
    RF2.7 & Al registrar o editar una tarea el usuario puede crear una categoría para la misma.\\
    \hline
    RF2.8 & El sistema permite al usuario ver las tareas organizadas por categorías. \\
    \hline
    RF2.9 & Al registrar o editar una tarea el usuario puede elegir una categoría existente para clasificarla.\\
    \hline
    RF2.10 & Las tareas con plazo se mostrarán como una etiqueta sobre la vista de calendario en los días especificados.\\
    \hline
    RF2.11 & El sistema permite al usuario marcar tareas como completadas. \\
    \hline
    RF2.12 & Una tarea que haya sido marcada como completada, se vera en la vista de calendario con una línea horizontal sobre el título de la tarea.\\
    \hline
    \multicolumn{2}{|c|}{\textbf{3. Recomendación de recursos para consulta}} \\
    \hline
    RF3.1 & Al agregar una tarea, el sistema genera recomendaciones de consulta (enlaces a artículos, video o páginas web).  \\
    \hline
    RF3.2 & El sistema genera recomendaciones con base al titulo de la tarea y al grado académico que el usuario ingresó al crear la cuenta.\\   
    \hline
    RF3.3 & El usuario puede guardar las sugerencias de consulta, haciendo que aparezcan cada vez que abre el detalle de la tarea.\\
    \hline
    RF3.4 & El sistema permite al usuario indicar si una recomendación de consulta no fue útil para él, el sistema tomará esta opinicón en
    cuenta para próximas recomendaciones. \\
    \hline
    \end{tabular}
    \caption{Requisitos funcionales consolidados para el sistema de gestión de tareas académicas y profesionales}
    \label{tab:requisitos-funcionales-2}
\end{table}
%--------------
\pagebreak

\section{Requerimientos No Funcionales}
\begin{table}[H]
    \centering
    \begin{tabular}{|p{3cm}|p{12cm}|}
    \hline
    \multicolumn{2}{|c|}{\textbf{Requerimientos No Funcionales (RNF)}} \\
    \hline
    \textbf{Categoría} & \textbf{Descripción} \\
    \hline
    Rendimiento & \textbf{RNF1} Las operaciones de interacción básica en la interfaz (abrir modal, guardar/editar tarea, navegar entre vistas) deberán responder en menos de \textbf{300 ms} en el 95\% de las peticiones bajo carga nominal. \\ 
    \hline
    Disponibilidad & \textbf{RNF2} El sistema deberá garantizar una disponibilidad mínima del \textbf{99.5\%} mensualmente (excluyendo ventanas de mantenimiento programado). \\ 
    \hline
    Seguridad & \textbf{RNF3} Autenticación segura (contraseñas almacenadas con hashing). \newline
                \textbf{RNF3} Todas las comunicaciones entre cliente y servidor deben estar cifradas mediante protocolo HTTPS. \newline
                \textbf{RNF3} Los usuarios deben autenticarse mediante credenciales seguras y el sistema debe implementar bloqueo tras 3 intentos fallidos. \\ 
    \hline
    Usabilidad & \textbf{RNF4} Interfaz intuitiva y consistente; realizar las principales acciones (crear tarea, marcar completada, editar) en \textbf{máximo 3 clics}. \newline
                \textbf{RNF4} Soporte móvil y escritorio. \\ 
    \hline
    Compatibilidad & \textbf{RNF5} Funcionamiento garantizado en navegadores modernos: Chrome, Firefox, Edge y Safari; versión adaptada para navegadores móviles equivalentes. \\ 
    \hline
    Mantenibilidad & \textbf{RNF6} Código modular y documentado.\\ 
    \hline
    Respaldo y Rendimiento & \textbf{RNF7} Copias de seguridad automáticas de la base de datos con retención mínima de \textbf{30 días}. \newline
                             \textbf{RNF7} El sistema debe soportar al menos 100 usuarios concurrentes sin degradación del rendimiento.\\ 
    \hline
    Recomendaciones & \textbf{RNF8} Las recomendaciones generadas al crear una tarea deberán presentarse en pantalla en menos de \textbf{5 segundos}.\\ 
    \hline
    Legal y Ética & \textbf{RNF9} Incluir términos y condiciones y política de privacidad. \newline
                   \textbf{RNF9} Las recomendaciones deben evitar contenidos que violen derechos de autor o políticas de contenido inapropiado. \\ 
    \hline
    \end{tabular}
    \caption{Requerimientos no funcionales principales}
    \label{tab:nfr}
\end{table}

\pagebreak
\section{Historias de Usuario}
\begin{table}[H]
    \centering
    \begin{tabular}{|p{15cm}|}
    \hline
    \textbf{Historia: } Registro de tareas\\
    \hline
    \textbf{Descripción: }Como profesional, necesito asignar fechas límite y estimar el tiempo que me tomará cada tarea para planificar mejor mi agenda y 
    cumplir con mis compromisos.\\
    \hline
    \textbf{Resultado: } Tareas programadas con fechas límite y tiempo estimado.\\
    \hline
    \textbf{Criterio de aceptación: } El usuario puede establecer fecha límite y duración estimada, y el sistema muestra indicadores visuales para 
    tareas próximas a vencer o vencidas.\\
    \hline
    \end{tabular}
    \caption{Programación de tareas con fechas}
    \label{user-story1}
\end{table}

\begin{table}[H]
    \centering
    \begin{tabular}{|p{15cm}|}
    \hline
    \textbf{Historia: } Modificación de tareas existentes\\
    \hline
    \textbf{Descripción: }Como usuario, necesito editar las tareas existentes para actualizar su información, cambiar fechas límite 
    o modificar descripciones según cambien mis necesidades.\\
    \hline
    \textbf{Resultado: }Tarea modificada con la información actualizada.\\
    \hline
    \textbf{Criterio de aceptación: } El usuario puede editar cualquier campo de sus tareas existentes, guardar los cambios y ver la 
    información actualizada inmediatamente en todas las vistas.\\
    \hline
    \end{tabular}
    \caption{Modificación de tareas existentes}
    \label{user-story2}
\end{table}

\begin{table}[H]
    \centering
    \begin{tabular}{|p{15cm}|}
    \hline
    \textbf{Historia: } Eliminación de tareas\\
    \hline
    \textbf{Descripción: }Como usuario, necesito eliminar tareas que ya no son relevantes o fueron 
    creadas por error para mantener mi lista organizada.\\
    \hline
    \textbf{Resultado: } Tarea eliminada permanentemente del sistema.\\
    \hline
    \textbf{Criterio de aceptación: }El usuario puede eliminar una tarea mediante confirmación, y esta 
    desaparece de todas las vistas sin posibilidad de recuperación.\\
    \hline
    \end{tabular}
    \caption{Eliminación de tareas}
    \label{user-story3}
\end{table}

\begin{table}[H]
    \centering
    \begin{tabular}{|p{15cm}|}
    \hline
    \textbf{Historia: } Visualización mensual de tareas\\
    \hline
    \textbf{Descripción: }Como usuario con múltiples tareas, necesito ver todas mis tareas distribuidas en un calendario mensual 
    para tener una perspectiva clara de mi carga de trabajo.\\
    \hline
    \textbf{Resultado: } Vista de calendario mostrando todas las tareas del mes.\\
    \hline
    \textbf{Criterio de aceptación: }El usuario puede acceder a una vista de calendario que muestra las tareas en sus fechas 
    correspondientes, con posibilidad de hacer clic para ver detalles o editar.\\
    \hline
    \end{tabular}
    \caption{Vista de calendario mensual}
    \label{user-story3}
\end{table}

\begin{table}[H]
    \centering
    \begin{tabular}{|p{15cm}|}
    \hline
    \textbf{Historia: } Creación de categorías personalizadas\\
    \hline
    \textbf{Descripción: }Como usuario, necesito crear categorías personalizadas para clasificar
     mis tareas según mis propios criterios de organización (trabajo, personal, estudio, materias, etc.).\\
    \hline
    \textbf{Resultado: } Nuevas categorías creadas y disponibles para asignar a tareas.\\
    \hline
    \textbf{Criterio de aceptación: } El usuario puede crear y asignar categorías con nombre único y color }
    distintivo, y estas aparecen disponibles al crear o editar tareas.\\
    \hline
    \end{tabular}
    \caption{Creación de categorías personalizadas}
    \label{user-story4}
\end{table}

\begin{table}[H]
    \centering
    \begin{tabular}{|p{15cm}|}
    \hline
    \textbf{Historia: } Visualización selectiva de tareas\\
    \hline
    \textbf{Descripción: }Como usuario con muchas tareas, necesito filtrar mi lista por categorías 
    específicas para enfocarme solo en las tareas de un tipo particular.\\
    \hline
    \textbf{Resultado: } Vista de tareas filtrada según la categoría seleccionada.\\
    \hline
    \textbf{Criterio de aceptación: } El usuario puede seleccionar una o múltiples categorías y ver solo las tareas 
    correspondientes, con opción de limpiar filtros para ver todas las tareas.\\
    \hline
    \end{tabular}
    \caption{Filtrado por categorías}
    \label{user-story5}
\end{table}

\begin{table}[H]
    \centering
    \begin{tabular}{|p{15cm}|}
    \hline
    \textbf{Historia: } Información detallada de tareas\\
    \hline
    \textbf{Descripción: }Como usuario, necesito agregar notas, instrucciones detalladas o comentarios adicionales a
     mis tareas para recordar información importante o pasos específicos.\\
    \hline
    \textbf{Resultado: } Tareas con información adicional almacenada.\\
    \hline
    \textbf{Criterio de aceptación: } El usuario puede agregar texto libre en un campo de notas que se guarda 
    con la tarea y es visible al editarla o revisarla.\\
    \hline
    \end{tabular}
    \caption{Notas adicionales en tareas}
    \label{user-story6}
\end{table}

\begin{table}[H]
    \centering
    \begin{tabular}{|p{15cm}|}
    \hline
    \textbf{Historia: }Sugerencias automáticas de recursos académicos\\
    \hline
    \textbf{Descripción: }Como estudiante o profesional, necesito recibir recomendaciones automáticas de artículos, videos o 
    recursos web relacionados con mis tareas para apoyar mi investigación y aprendizaje.\\
    \hline
    \textbf{Resultado: }Lista de recursos recomendados generados automáticamente.\\
    \hline
    \textbf{Criterio de aceptación: } Al crear o ver una tarea, el sistema genera y muestra 3-5 recursos relevantes
     (enlaces, artículos, videos) basados en el título y descripción de la tarea.\\
    \hline
    \end{tabular}
    \caption{Recomendaciones de recursos de consulta}
    \label{user-story7}
\end{table}

\begin{table}[H]
    \centering
    \begin{tabular}{|p{15cm}|}
    \hline
    \textbf{Historia: } Recomendaciones personalizadas por perfil\\
    \hline
    \textbf{Descripción: }Como usuario con un perfil académico específico, necesito recibir recomendaciones de 
    recursos y organización adaptadas a mi nivel de estudios y área de especialización.\\
    \hline
    \textbf{Resultado: } Recomendaciones personalizadas según el perfil del usuario.\\
    \hline
    \textbf{Criterio de aceptación: } El sistema utiliza la información del perfil académico del usuario (grado de estudios, 
    área de especialización) para generar recomendaciones más precisas y apropiadas para su nivel.\\
    \hline
    \end{tabular}
    \caption{Recomendaciones personalizadas por perfil}
    \label{user-story8}
\end{table}

\begin{table}[H]
    \centering
    \begin{tabular}{|p{15cm}|}
    \hline
    \textbf{Historia: } Almacenamiento de recursos útiles\\
    \hline
    \textbf{Descripción: }Como usuario, necesito poder guardar las recomendaciones de consulta que considero útiles para poder 
    acceder a ellas fácilmente cada vez que abra el detalle de la tarea sin tener que buscarlas nuevamente.\\
    \hline
    \textbf{Resultado: } Recomendaciones seleccionadas guardadas y visibles en el detalle de la tarea.\\
    \hline
    \textbf{Criterio de aceptación: }El usuario puede marcar recomendaciones como favoritas mediante un botón o ícono, estas se guardan 
    permanentemente y aparecen destacadas cada vez que abre el detalle de la tarea.\\
    \hline
    \end{tabular}
    \caption{Guardado de recomendaciones favoritas}
    \label{user-story9}
\end{table}

\begin{table}[H]
    \centering
    \begin{tabular}{|p{15cm}|}
    \hline
    \textbf{Historia: } Feedback de recomendaciones para mejora del sistema\\
    \hline
    \textbf{Descripción: }CComo usuario, necesito poder indicar si una recomendación de consulta fue útil o
     no para que el sistema aprenda de mis preferencias y mejore las futuras sugerencias.\\
    \hline
    \textbf{Resultado: } Feedback del usuario registrado para mejorar el algoritmo de recomendaciones.\\
    \hline
    \textbf{Criterio de aceptación: } El usuario puede marcar una recomendación como "útil" o "no útil" mediante botones simples, 
    el sistema almacena esta retroalimentación y la utiliza para personalizar futuras recomendaciones del usuario.\\
    \hline
    \end{tabular}
    \caption{Feedback de recomendaciones}
    \label{user-story10}
\end{table}
%%%%%%%%%% BIBLIOGRAFÍA %%%%%%%%%%%%%%%%%%%%%
\pagebreak
\printbibliography[heading=bibintoc]

\end{document}
